\part{Building it from nothing}

\noindent\textbf{\large What this part covers}

\medskip
\noindent
The order the pieces have to be built in, and why that order is forced rather
than chosen. If you set out to rebuild this project by hand, this is the
sequence --- and skipping ahead in it produces work you have to throw away.

\chapter{The order, and why it is not arbitrary}
\label{ch:buildorder}

\section{The principle}

Build \textbf{inward-out}: the things that cannot change first, the things that
change constantly last.

\begin{center}
\begin{tabular}{@{}rp{5.0cm}p{6.6cm}@{}}
\toprule
& \textbf{Stage} & \textbf{Why here} \\
\midrule
1 & The database schema & Everything else refers to it. Getting it wrong later means migrations, data rewrites and edits everywhere. \\
2 & Enums and models & The vocabulary. Once \fpath{IncidentStatus} exists, every later file can use it. \\
3 & Authentication & Nothing else can be built until you can say who is asking. \\
4 & Permissions & Every endpoint needs to ask ``may they?'' before it exists. \\
5 & Services (the logic) & The actual rules, testable with no HTTP involved. \\
6 & Controllers and routes & A thin shell over the above. \\
7 & The realtime service & Needs events to exist before it can forward them. \\
8 & The frontend & Needs an API to talk to. \\
9 & Docker and nginx & Wires together things that already work individually. \\
10 & CI and deployment & Automates a build you can already do by hand. \\
\bottomrule
\end{tabular}
\end{center}

\note{The temptation is always to build the screen first, because a screen is
visible and a migration is not. It is a false economy: a screen built against an
API that does not exist is built against your \emph{guess} at that API, and the
guess is usually wrong in ways that are invisible until integration.}

\section{Stage 1: the schema}

Eight migrations, and their numeric prefixes are the order they run in.

\begin{lstlisting}
2026_01_01_000100_...   organizations and users
2026_01_01_000200_...   membership: who belongs to which organisation
2026_01_01_000300_create_incidents_table.php
2026_01_01_000400_create_incident_activity_tables.php
...
\end{lstlisting}

\warnbox{The order is forced by foreign keys. \fpath{incidents} has an
\fpath{organization\_id} that must point at a real organisation, so
\fpath{organizations} has to exist first. Get the numbering wrong and the
migration fails --- immediately and clearly, which is the good case. The bad
case is a schema that lets a row point at nothing.}

Design decisions that had to be made at this stage, because changing them later
is expensive:

\begin{itemize}[leftmargin=1.4cm]
\item \fpath{incident\_events} has no \fpath{updated\_at} and no
  \fpath{deleted\_at}. There is nowhere for an edit to be recorded, because
  edits are not permitted. The schema states the intent before any code enforces
  it.
\item \fpath{time\_to\_acknowledge\_seconds} lives on the incident row rather
  than being computed from the timeline. Chapter~\ref{ch:screens} explains why
  that is safe.
\item Idempotency keys have a unique index. That index \emph{is} the mechanism
  --- see Chapter~\ref{ch:bugs-backend}.
\end{itemize}

\section{Stage 2: enums before models}

Write \fpath{IncidentStatus} and \fpath{IncidentSeverity} before
\fpath{Incident}. The model then declares its columns \emph{as} those types, and
from that moment on it is impossible to assign an invalid status anywhere in the
codebase --- PHP refuses.

\note{This is why the transition rule lives on the enum rather than in a service.
The enum is the one thing every other layer already depends on, so a rule placed
there is unavoidable. A rule in a service can be bypassed by not calling that
service.}

\section{Stage 3: authentication, and the shape of the key}

Authentication had to be built before anything it protects. The decision made
here that shaped the rest of the system: \textbf{RS256, not HS256}.

\begin{center}
\begin{tabular}{@{}>{\raggedright\arraybackslash}p{2.6cm}>{\raggedright\arraybackslash}p{4.6cm}p{6.4cm}@{}}
\toprule
& \textbf{HS256} & \textbf{RS256} \\
\midrule
Key & One shared secret & A private key and a public key \\
To verify & You need the secret & You need only the public key \\
Consequence & Anyone who can verify can also forge & Verifying and signing are separate powers \\
\bottomrule
\end{tabular}
\end{center}

The realtime service must verify tokens. With HS256 it would need the signing
secret, and a compromise of the most exposed service would hand over the ability
to mint tokens for anyone. With RS256 it holds only the public half.

\warnbox{This choice creates an operational requirement that catches people
during deployment: \emph{all three} PHP containers must have the same private
key. In Docker Compose they share a volume, so it happens automatically. On a
platform that gives each container its own filesystem, each would generate a
different key and tokens minted by one would fail verification in another. The
symptom is intermittent, inexplicable ``invalid signature'' errors.
Chapter~\ref{ch:bugs-deploy}.}

\section{Stage 5: the logic, before any HTTP}

The most important file in the project is \fpath{IncidentService}, and it knows
nothing about HTTP. It does not see requests or return responses. That is what
makes the rules testable without a browser, a server, or a network.

The pattern it follows appears throughout:

\begin{enumerate}[leftmargin=1.5cm]
\item Check the rules.
\item Open a transaction.
\item Write everything --- the change, the timeline entry, the notification rows.
\item Commit.
\item \emph{Then} queue jobs and publish events.
\end{enumerate}

\note{Steps 4 and 5 are in that order for a reason given in
Chapter~\ref{ch:arch}, and the project has a dedicated class ---
\fpath{PendingEffects} --- whose entire job is to hold the side effects until
the transaction has committed. When a codebase grows a class to enforce an
ordering, it is usually because getting it wrong once was expensive.}

\section{Stage 9: Docker last, not first}

Docker is written when the pieces already work. It is a way of \emph{running}
software, not a way of writing it.

\begin{center}
\begin{tabular}{@{}>{\raggedright\arraybackslash}p{4.0cm}p{9.6cm}@{}}
\toprule
\textbf{File} & \textbf{Written when} \\
\midrule
\fpath{api/Dockerfile} & The API runs and its tests pass. \\
\fpath{web/Dockerfile} & The frontend builds and talks to the API. \\
\fpath{docker-compose.yml} & Every service works individually. \\
\fpath{infra/nginx/} & There is more than one service to route between. \\
\bottomrule
\end{tabular}
\end{center}

\warnbox{Writing the compose file first is a common instinct and a trap: you end
up debugging container networking and application logic simultaneously, unable
to tell which layer a failure came from. Get each piece working on its own, then
wire them together --- so that when something breaks, only one thing changed.}

\section{Stage 10: what CI actually protects}

The pipeline runs the tests, the linters, a security audit of dependencies, a
container vulnerability scan, and then starts the entire stack and runs the
browser tests against it.

That last step is the expensive one and the one that has caught the most. Every
bug in Chapter~\ref{ch:bugs-infra} lived in the gaps \emph{between} services,
where no unit test could reach --- and several of them were found only because CI
insisted on starting the whole system from nothing, on a machine with no
leftover state.

\note{This is the deepest lesson in the book. Local development accumulates
state: containers that were already running, a database that was already seeded,
a key that was already generated. CI has none of that, every single time. A
system that only works on a machine with history is not a working system --- it
just has not been asked to prove itself yet.}
